\section{Political reanalysis: how many democracies really declined?}\label{sec:political}

We analyze trust in parliament (\texttt{trstprl}) and satisfaction with democracy
(\texttt{stfdem}) in ESS rounds 9--11 (2018--2024, the window whose integrated
files ship PSU/stratum identifiers), then extend every trajectory to the full
record (rounds 1--11, 2002--2024), asking, claim by claim, how many countries
survive.
The comparative literature has grown skeptical of a continent-wide trust decline
\citep{norris2017western,voeten2017people,wuttke2022europeans}; we add a
formal procedure that says so and a diagnosis of why the over-reading arises:
the dominant trend estimates \citep{valgardsson2025crisis} rest on
latent-trait models with no design-based uncertainty layer
\citep{taihusolt2024}, and latent-mood estimation
\citep{claassen2019estimating,claassen2020support,kolczynska2026trust} pools
countries into model-based credible intervals, where our bands give
model-free coverage at a stated claim; the two are complements.
Certification levels here are design-bootstrap simultaneous levels,
asymptotic in the within-country samples (Table~\ref{tab:targets}).

\paragraph{The counts collapse across claims.} For \texttt{trstprl} over rounds
9--11, a marginal pairwise reading flags $20$ of $30$
countries; an any-pair
design-aware band drops this to $12$; net first-to-last decline holds in $6$
(Austria, Belgium, Estonia, the United Kingdom, Greece, the Netherlands); and the
\emph{persistent, design-aware simultaneous} band over the preregistered low-trust
core certifies one, Greece, which also survives Bonferroni across countries.
\emph{Greece fielded only rounds 10 and 11 here}, so its trajectory is a single pair and the top claims coincide for it
(as for Czechia and Israel, uncertified). Among the $27$ two-pair
countries none certifies persistence, while the United Kingdom, Belgium, and
the Netherlands pass the \emph{plug-in} persistent claim and are demoted by
the design-aware layer alone. The headline is therefore two numbers: net
decline in six of thirty, and \emph{zero} of the twenty-seven countries
observed across multiple pairs certifying persistence. The severity analysis
(Section~\ref{sec:simulation}) explains why: the persistent claim has $80\%$ power
only against crisis-scale declines, so its non-certifications are weak evidence of
absence. For \texttt{stfdem} the pattern repeats ($23$ marginal, $14$ any-pair,
$8$ net, one persistent: Greece, again one pair, failing Bonferroni there).
The literature reads Greece as a crisis-era legitimacy collapse
\citep{armingeon2014democracy,torcal2014decline}, but those accounts concern
2009--2015 while the certified pair is 2020/21$\to$2023/24.

\paragraph{The long window: the secular question, asked properly.} Extending every
trajectory to its full record (rounds 1--11; median six pairs, up to ten;
rounds 1--8 weights-only in the integrated files; merging their sample-design
files changes no net or persistent count below --- Supplementary Material)
reframes the twenty-year narrative in both directions. Zero of the thirty-three countries with three or more usable rounds
certify persistent decline over the preregistered core across their full
record, plug-in included. Reversing the inequality certifies
\emph{recoveries}, and across the $1{,}173$ ordered contrasts of the trust series
certified movement is directionally balanced ($193$ declining against $212$
rising). Inside the same band, at the same $\alpha$, twenty-three of
thirty-three countries certify at least one decline and at least one
recovery on each outcome (twenty on both); the two directions share one
critical value, so no within-country correction is needed. The count across countries is a
conjunction of thirty-three coverage events and carries no familywise control
by itself; the closed-testing bound below supplies it. The certified record
therefore argues against describing European trust as a uniformly monotone
slide --- on the positive evidence of coexisting certified declines and
recoveries, not on the persistent rung's non-certifications, which its
severity makes weak evidence of absence.

Under Proposition~\ref{prop:transfer}---one band, every ordered contrast and
both directions at $\alpha=0.10$ jointly---net erosion over the
whole span certifies in eight countries (Cyprus, Spain, the United
Kingdom, Greece, Hungary, Israel, Italy, Ukraine;
Figure~\ref{fig:bounds}, Table~\ref{tab:magnitudes}).
Against the per-claim construction the joint statement costs France at the span
claim ($9\to8$) and six countries at the any-pair claim ($28\to22$); the
persistent zero holds under both constructions (Supplementary Material).

The across-country count itself can now carry a guarantee. Inverting each
joint band over its level yields a per-country certification $p$-value;
closed testing across the thirty-three \citep{goeman2011multiple} then
certifies, with $90\%$ simultaneous confidence, that \emph{at least six}
truly declined over their span, on each outcome, the six nameable at no
further cost (Supplementary Material). The Simes local tests behind that
bound need the countries' $p$-values independent or positively dependent
--- a design assumption (independent sampling mechanisms given the finite
populations), not a property of separate bootstrap streams; Bonferroni local
tests, which assume nothing about dependence, return the same six on both
outcomes, as does the rounds 1--8 design-file upgrade (Supplementary
Material). The wrong-unit ladder thus closes at
the top: threshold, wave, trajectory, country, count.

Under-detection shows in the ledger of contrasts: Hungary certifies its
$22$-year span while only $9$ of its $55$ ordered contrasts certify, and the
United Kingdom $11$ of $55$ --- erosion a reading built from adjacent waves
would mostly miss. Because every endpoint choice is certified at once, the
band also delivers the \emph{share of a country's ordered contrasts that
certify a decline}: a lower bound, comparable within a country but not across
(the shared critical value increases strongly with record length while the
denominator grows quadratically), so we do not rank on it. Share and
magnitude differ: Spain certifies $32$ of $55$ contrasts
yet clears its span by only $0.6$ CDF points, and Slovenia certifies $21$ of
$55$ without entering the eight at all, the endpoint sensitivity the joint
band exposes. Ukraine's certified span decline of at least $25$ points runs
from the post-Orange-Revolution wave of 2005 to a wartime sample: a certified
statement about those two waves.

\begin{table}[t]\centering
\caption{Certified net decline in parliamentary trust, 2002--2024, read off
each country's joint band (every ordered contrast, both directions,
$\alpha=0.10$). \emph{Decline $\ge$}: simultaneous lower bound over the
preregistered core, in CDF points. \emph{Contrasts}: ordered round pairs
certified declining, of all admissible; their share is the endpoint-free
erosion index. The certified set is stable under weights-only design-effect
inflation up to roughly $1.3$; above that Spain, the marginal member, drops
out (Supplementary Material).}\label{tab:magnitudes}
\begin{tabular}{lcccc}
\toprule
country & rounds & decline $\ge$ & contrasts & share \\
\midrule
Ukraine & 2--11 & 25.9 & 7/15 & 0.47 \\
Cyprus & 3--11 & \phantom{0}6.2 & 13/21 & 0.62 \\
United Kingdom & 1--11 & \phantom{0}6.1 & 11/55 & 0.20 \\
Greece & 1--11 & \phantom{0}3.1 & 8/15 & 0.53 \\
Israel & 1--11 & \phantom{0}2.6 & 10/28 & 0.36 \\
Italy & 1--11 & \phantom{0}1.6 & 5/15 & 0.33 \\
Hungary & 1--11 & \phantom{0}1.1 & 9/55 & 0.16 \\
Spain & 1--11 & \phantom{0}0.6 & 32/55 & 0.58 \\
\bottomrule
\end{tabular}
\end{table}
\begin{figure}[t]\centering
\includegraphics[width=.8\textwidth]{ess_net_bounds.pdf}
\caption{Each country's simultaneous lower bound on the net (first-to-last)
decline in parliamentary trust, 2002--2024, read off its joint band at
$\alpha=0.10$. Filled markers: the eight countries whose bound exceeds zero;
asterisks: the six named by the closed-testing prevalence bound. Negative
bounds say how far a country sits from certifying, not that it
improved.}\label{fig:bounds}
\end{figure}

The reclassified countries are the weak-signal, higher-noise cases---\emph{certified
by the marginal plug-in but inconclusive under the design-aware band}, never
``false positives,'' a term we reserve for simulation with known truth.

\paragraph{Two confounds this window cannot exclude.} Round-10 fieldwork
coincided with the pandemic's rally-round-the-flag movements, so ``persistent
decline through this window'' is a claim made against a rally; and nine
countries switched to self-completion at round 10, confounding wave-pair
shifts with mode effects, which no design bootstrap absorbs. A mode audit
from the data's own mode variable bounds it: the Greek pair is mode-constant
and net counts are unchanged when the switchers' confounded pairs are
excluded (Supplementary Material, with a preregistered age-group analysis
that finds no pooling artifact).

\paragraph{A second, global target: Foa--Mounk deconsolidation.} We preregistered a
reanalysis of the ``democratic deconsolidation'' thesis
\citep{foa2016democratic,foa2017signs} on the WVS/EVS trends (1981--2022,
$59$--$77$ countries with $\ge2$ waves), recoding the five-item battery so that
deconsolidation is a persistent decline over each scale's preregistered core.
The harmonized trend file provides no design metadata, so this is an external
substantive stress test of the thesis, not a second validation of the method;
the ESS carries the design-based validity story. The verdict is calibration,
not denial, but the surviving phenomenon is not the one the thesis was about.
A \emph{marginal wave-pair reading of the same data}, constructed to
instantiate the bottom rung, flags $36$--$64$ countries per item against
$6$--$17$ (8--24\%) at the persistent rung. That $2.6$--$6.5\times$ ratio compounds two
changes: holding the inference layer fixed, the rung alone accounts for
$1.7$--$4.8\times$ (plug-in) or $1.9$--$4.8\times$ (design-aware), and
propagating design uncertainty at fixed rung accounts for the remainder.
Three caveats. The trend file carries no PSU/stratum identifiers, so the
weights-only bootstrap omits within-PSU dependence (variance-inflating) and
stratification (variance-reducing); where the former dominates, as is usual,
certified counts are \emph{inflated}: at variance $\times1.5$ and
$\times2.0$ the certified core holds at $13$ and $12$ countries and the gap
widens to $3.0$--$8.3\times$, so the reported gap is then a lower bound. The
per-item lists
carry no across-country familywise control, so we aggregate to multi-item
co-certification before interpreting geography. And the certification sets
sit below the $K\ge94$ bound, so every WVS band uses the full-$\alpha$
anchors of Theorem~\ref{thm:safe}.

\paragraph{The certified core: where the surviving deconsolidation lives.}
Aggregating per-item certifications yields a positive characterization
(Figure~\ref{fig:core}): of $38$ countries certified on $\ge1$ item, $13$
form a certified core on $\ge2$, headed by Albania on four. The core is
post-communist and Arab-Spring states, not the mature democracies the thesis
concerned,
formalizing with per-country error control the decline-from-euphoric-baselines
pattern the post-communist literature described qualitatively
\citep{mishler1997trajectories,inglehart2003honeymoon}. Three qualifications
travel with it. Five core members are electoral autocracies in the V-Dem
classification, where falling stated support reflects autocratic legitimation
rather than deconsolidation in the consolidated-democracy sense. The West enters
the core exactly twice (Finland and Switzerland, on the strongman/army pair,
which we flag as a probable administration artifact): no consolidated Western
democracy certifies on support for a democratic system, consistent with
\citet{wuttke2022europeans}, though not that the West is absent. And the
geography is partly window-contingent: recertifying on fieldwork ending by
2017 keeps eight of the thirteen but drops all three Arab-Spring cases, so
the post-communist component is window-stable and the Arab-Spring component
rests on WVS wave~7. Certified attitudes do not forecast institutions (core
membership predicts subsequent V-Dem electoral-democracy decline no better
than the marginal flag, Fisher $p=0.89$, a comparison too small to establish
the decoupling recent panel work expects,
\citealp{ouattara2023distrusting,vandermeer2025pushing}): certification is
descriptive error control, not a forecast (item-level caveats in the
Supplementary Material).

\begin{figure}[t]\centering
\includegraphics[width=.95\textwidth]{certified_core.pdf}
\caption{The thirteen-country certified core, by evidence strength (WVS/EVS
1981--2022; $K=59$--$77$ per item). Circle area is the persistent band's
simultaneous lower bound on the per-pair decline, in CDF points; open circles
mark items that do not certify. The full thirty-eight-country matrix is in the
Supplementary Material.}\label{fig:core}
\end{figure}
